% ======================================================================
% SECTION 10: ROBOTIC EXECUTION AND PHYSICAL AI GATEWAY
% File: sections/robotic_execution.tex
% ======================================================================

The robotic execution gateway is the component that distinguishes a Physical AI sponsor from
a conventional software-based clinical trial management system. While all preceding sections
describe agent functions that could, in principle, operate in a purely digital trial context,
this section addresses the unique capability of the autonomous sponsor to issue bounded,
safety-checked instructions to physical robotic systems that perform clinical procedures on
cancer patients. The robot execution gateway translates sponsor-level decisions into robotic
task orders, evaluates safety prerequisites before every procedure, manages procedure state
through a finite state machine, and handles emergency conditions with mandatory human override.

\subsection{Robot Execution Gateway}

The robot execution gateway (robot\_execution\_gateway) provides the central interface between
the autonomous sponsor agent network and the physical robotic systems operating at trial
sites. Unlike other agents that replace existing human roles, the robot execution gateway
has no traditional equivalent in pharmaceutical sponsor organizations, as traditional
sponsors do not directly control clinical procedure execution. This agent represents the
fundamental architectural extension required for Physical AI trial sponsorship.

Task-order generation translates sponsor agent decisions into bounded robotic instructions.
When the clinical accountability agent authorizes a treatment, the study orchestrator
schedules the procedure, and the site gateway confirms site readiness, the robot execution
gateway generates a formal task order that specifies the procedure type, target patient,
required robotic system, operating parameters (force limits, workspace boundaries, speed
constraints), success criteria, and timeout conditions. Each task order is a self-contained
instruction set that the robotic system executes within defined safety envelopes.

Site capability matching determines whether a specific trial site has the qualified robotic
infrastructure to perform a requested procedure. The gateway maintains a real-time registry
of all qualified robots at all trial sites, their current operational status, maintenance
certification currency, and procedure qualification records. When a procedure is requested,
the gateway matches the procedure requirements against available capabilities, selecting the
optimal robot based on qualification status, current availability, and historical performance.

Safety gate evaluation implements a multi-step approval process before every robotic
procedure. Gate 1 verifies patient eligibility and consent for the specific procedure.
Gate 2 confirms robotic system readiness (calibration current, no active alarms, required
instruments available). Gate 3 validates the procedure parameters against protocol-defined
safety envelopes. Gate 4 requires supervising clinician confirmation for Tier 2 and Tier 3
procedures (per the adapted 21 CFR Part 50 three-tier Physical AI classification system
\cite{cfr50-adapt}). All four gates must pass before the gateway authorizes procedure
initiation. Gate failure at any point halts the procedure sequence and routes the issue to
the appropriate resolution path.

Procedure-state management implements a finite state machine with six states: Idle (no active
procedure), Pre-operative (pre-procedure verification and patient positioning), Active
(robotic procedure in progress with continuous monitoring), Post-operative (post-procedure
assessment and documentation), Recovery (patient recovery monitoring), and Complete (procedure
concluded with all documentation finalized). Transitions between states are governed by
defined criteria and produce audit trail entries documenting the state change, triggering
conditions, and responsible agent.

Emergency stop handling provides the highest-priority response to conditions that may
threaten patient safety during a robotic procedure. The gateway monitors all robotic
telemetry streams for emergency conditions: force limits exceeded, position deviation beyond
safe envelope, communication loss with the robotic system, patient vital sign alarm, or
manual emergency stop activation by the supervising clinician. Upon detection of any
emergency condition, the gateway issues an immediate halt command, preserves the current
robotic state for forensic analysis, activates the clinical emergency protocol, and generates
a structured incident report. Autonomous procedure resumption after an emergency stop is
never permitted; resumption requires explicit human authorization through a mandatory
safety review.

\subsection{Robotic Capability Registry}

The robotic capability registry maintains the comprehensive inventory of all qualified
robotic systems across the trial network. Table~\ref{tab:robot-registry} presents the
registry of the ten robotic system categories evaluated in the national platform
\cite{national-platform-paper}, based on the 24-hour simulation involving 168 patients
and 29 robot instances \cite{main-repo}.

\begin{longtable}{L{2.0cm} L{1.2cm} L{1.3cm} L{1.2cm} L{2.5cm} L{3.0cm}}
\caption{Robotic System Registry}
\label{tab:robot-registry} \\
\toprule
\textbf{Robot Category} & \textbf{Count} & \textbf{Type} & \textbf{USL Score} & \textbf{Trial Function} & \textbf{Safety Classification} \\
\midrule
\endfirsthead
\toprule
\textbf{Robot Category} & \textbf{Count} & \textbf{Type} & \textbf{USL Score} & \textbf{Trial Function} & \textbf{Safety Classification} \\
\midrule
\endhead
\midrule
\multicolumn{6}{r}{\textit{Continued on next page}} \\
\endfoot
\bottomrule
\endlastfoot
Surgical Robots & 3 & Surgical & 87.5 & Mediastinal tumor resection, thoracoscopic procedures & Tier 3: mandatory human gate \\
Cobots & 4 & Collaborative & 88.75 & Soft-tissue sarcoma biopsy, pharmacy dosing, lab automation & Tier 2: clinician confirmation \\
RT Positioning & 3 & Positioning & 82.5 & Brain tumor radiation patient alignment and immobilization & Tier 2: clinician confirmation \\
Needle Placement & 2 & Interventional & 85.0 & Parotid tumor percutaneous biopsy guidance & Tier 3: mandatory human gate \\
Social Companions & 5 & Companion & 90.0 & Pediatric leukemia anxiety reduction, family engagement & Tier 1: autonomous operation \\
Humanoids & 3 & Humanoid & 75.0 & Pediatric osteosarcoma mobility assistance, rehabilitation & Tier 2: clinician confirmation \\
RT Motion Tracking & 3 & Tracking & 86.25 & Lung tumor respiratory motion compensation during treatment & Tier 2: clinician confirmation \\
Imaging Assistants & 4 & Diagnostic & 83.75 & Liver tumor imaging positioning, contrast timing & Tier 1: autonomous with monitoring \\
Steerable Needles & 2 & Interventional & 85.0 & Liver tumor ablation needle guidance & Tier 3: mandatory human gate \\
Rehab Exoskeletons & 3 & Rehabilitation & 81.25 & Femur osteosarcoma post-surgical mobility recovery & Tier 2: clinician confirmation \\
\end{longtable}

Table~\ref{tab:procedure-safety-gates} details the procedure safety gate protocol.

\begin{longtable}{L{1.5cm} L{3.0cm} L{3.0cm} L{1.5cm} L{2.8cm}}
\caption{Procedure Safety Gate Protocol}
\label{tab:procedure-safety-gates} \\
\toprule
\textbf{Gate} & \textbf{Trigger} & \textbf{Required Approvals} & \textbf{Timeout} & \textbf{Fallback Action} \\
\midrule
\endfirsthead
\toprule
\textbf{Gate} & \textbf{Trigger} & \textbf{Required Approvals} & \textbf{Timeout} & \textbf{Fallback Action} \\
\midrule
\endhead
\midrule
\multicolumn{5}{r}{\textit{Continued on next page}} \\
\endfoot
\bottomrule
\endlastfoot
Gate 1: Patient verification & Procedure scheduled and patient present & Patient consent confirmed, identity verified by site\_gateway & 5 minutes & Procedure postponed, reschedule initiated \\
Gate 2: System readiness & Gate 1 passed & Robot calibration verified, instrument check, no active alarms & 10 minutes & Maintenance request, backup robot if available \\
Gate 3: Parameter validation & Gate 2 passed & Procedure parameters within safety envelopes, digital twin verification & 5 minutes & Parameter review, clinical accountability agent consultation \\
Gate 4: Clinical authorization & Gate 3 passed, Tier 2/3 procedures & Supervising clinician provides electronic authorization & 30 seconds for Tier 3 & Mandatory hold, no autonomous bypass \\
\end{longtable}

\subsection{Surgical Robot Operations}

Surgical robot operations center on the da Vinci Xi surgical system integration for
oncology procedures. The 24-hour simulation demonstrated robotic-assisted lobectomy
procedures in the surgical suites, with the da Vinci Xi achieving a Unification Standard
Level score of 87.5 reflecting strong performance across simulation switching, AI
integration, cross-robot sharing, and clinical trial collaboration dimensions
\cite{national-platform-paper}. The patient journey paper documented a 168-minute
robotic-assisted lobectomy achieving R0 resection with negative margins
\cite{patient-journey-paper}.

Pre-operative digital twin simulation provides a virtual rehearsal of the planned
procedure using patient-specific anatomical data. The digital twin reconstructs the
surgical field from pre-operative imaging (CT, MRI, PET), simulates instrument
trajectories and tissue interactions, and identifies potential complications before
the physical procedure begins. The robot execution gateway validates that the
digital twin simulation completes successfully before authorizing the physical procedure.

Intraoperative telemetry monitoring provides continuous safety oversight during active
procedures. The gateway monitors force and torque at instrument tips, instrument position
relative to defined safe zones, insufflation pressure, and system status indicators. When
any monitored parameter approaches safety limits, the gateway issues graduated warnings:
advisory (parameter approaching limit), caution (parameter near limit with trend toward
exceedance), and halt (parameter exceeded, procedure paused).

\subsection{Collaborative Robot Operations}

Collaborative robot (cobot) operations leverage the Franka Emika platform for laboratory
automation, pharmacy dosing, and sample processing workflows. The Franka Emika achieves
the highest USL score (88.75) among all evaluated robot categories, reflecting its advanced
collaborative capabilities and safety features \cite{national-platform-paper}.

Laboratory automation workflows encompass sample handling from collection through
processing: specimen accessioning with barcode verification, centrifugation with
g-force and duration control, aliquoting with volume verification, and storage in
temperature-controlled repositories. Each workflow step generates provenance records
that trace sample identity and handling conditions throughout the processing chain.

Pharmacy dosing workflows automate the preparation of weight-based chemotherapy doses.
The Franka Emika cobot receives a dose order from the supply agent (specifying patient
weight, prescribed dose, drug formulation, and diluent), calculates the required
preparation volume, selects the correct vial and diluent from the robotic pharmacy
inventory, performs reconstitution and dilution under sterile conditions, verifies the
final concentration through gravimetric measurement, applies a patient-specific label
(verified by OCR), and stages the prepared dose for pharmacist verification before
patient administration. The entire preparation sequence is recorded in the procedure
provenance system, creating a complete audit trail from drug order through patient
administration.

The cobot safety features include force-limited operation (the robot halts immediately if
contact force exceeds 50 Newtons), workspace boundary enforcement (the robot cannot operate
outside its defined workspace envelope), and dual-channel monitoring (independent safety
processors monitor the robot's primary control system). These safety features enable the
Franka Emika to operate in close proximity to pharmacy staff without physical barriers,
supporting the collaborative workflow model where robotic preparation is verified by human
pharmacists before patient administration.

Additional robotic categories operating within the trial site include radiation therapy
positioning robots (USL 82.5) that align and immobilize patients for brain tumor radiation
treatment, needle placement systems (USL 85.0) that guide percutaneous biopsies of parotid
and liver tumors, social companion robots (USL 90.0, the highest among all categories) that
reduce anxiety in pediatric leukemia patients, humanoid robots (USL 75.0) that assist with
pediatric osteosarcoma rehabilitation, radiation therapy motion tracking robots (USL 86.25)
that compensate for respiratory motion during lung tumor treatment, imaging assistant robots
(USL 83.75) that optimize patient positioning and contrast timing for liver tumor imaging,
steerable needle robots (USL 85.0) that guide ablation needles into liver tumors, and
rehabilitation exoskeletons (USL 81.25) that support post-surgical mobility recovery for
femur osteosarcoma patients.

Each robotic category requires a distinct integration profile within the robot execution
gateway. Surgical robots require the most extensive safety gate protocol (Tier 3: mandatory
human authorization for every procedure), while social companion robots operate with the
most autonomous profile (Tier 1: autonomous operation with logging). The gateway maintains
category-specific configuration profiles that define the appropriate safety gates, telemetry
parameters, emergency stop protocols, and documentation requirements for each robot type.
This configurable architecture enables the same gateway infrastructure to manage the full
spectrum of robotic systems, from high-autonomy companion robots to high-oversight surgical
systems.

\subsection{Device Telemetry and Procedure Provenance}

Device telemetry ingestion collects real-time sensor data from all robotic systems during
clinical procedures. The ingestion system handles diverse telemetry formats: ROS
(Robot Operating System) topics from research platforms, proprietary APIs from commercial
surgical systems, and standardized interfaces where available. All telemetry data is
timestamped to microsecond precision and stored in append-only databases that prevent
retrospective modification.

Procedure provenance reconstruction creates a complete timeline of every robotic procedure
from pre-operative setup through post-operative documentation. The timeline integrates
robotic telemetry, patient vital signs, clinical observations, safety gate events, and
operator interactions into a single chronological record. This record satisfies the
procedure provenance requirements defined in the adapted 21 CFR Part 312 \cite{cfr312-adapt}
and provides the evidence base for regulatory submission of robotic procedure outcomes.
The provenance chain integrates with the trust layer (\S\ref{sec:trust}) for immutable
storage and cryptographic integrity verification.

The device telemetry system also supports predictive maintenance for robotic systems. By
analyzing telemetry trends (increasing force variability, drift in position accuracy,
changes in response latency), the system can predict when a robotic system will require
maintenance before a failure occurs. Predictive maintenance scheduling is coordinated with
the site gateway to minimize impact on patient scheduling: maintenance is scheduled during
low-utilization periods when no procedures are planned. This approach maintains the 99.7\%
facility uptime demonstrated in the 24-hour simulation \cite{main-repo} while ensuring
that all robotic systems remain within their validated operating parameters.

The robot execution gateway maintains a comprehensive performance database that tracks
every procedure performed by every robot across all trial sites. This database enables
cross-site comparison of robotic system performance, identification of best practices for
procedure optimization, and detection of site-level performance variations that may
indicate training needs or equipment issues. The performance database feeds back into the
PSL and USL scoring systems, enabling dynamic updates to site and robot readiness scores
based on actual operational performance rather than static qualification assessments.
